\section{Introduction}

\begin{figure*}[!t]
\centering
\includegraphics[width=\textwidth]{fig1_overview.pdf}
\caption{The registered gate and certify-then-hedge on the same input, one fresh-seed episode
chosen post hoc to show the mechanism: of the 32 lead-time episodes, it is the gate's third-worst
and the hedge's tenth-lowest. \textbf{(a)}~At period~13 the alert yields a \shockspec{} from
Gemini~3.8 Flash (six of its nine fields shown) naming a lead-time shift of \emph{medium}
magnitude; the true shift is one period, from period~14. \textbf{(b)}~The gate tests the stated
hypothesis alone and, once $E_{j,t} \ge 1/\alpha_j$, from period~24 (eleven periods after the
alert) switches all at once to the configuration compiled from the stated magnitude, two extra
periods. \textbf{(c)}~Certify-then-hedge keeps the fixed set $\mathcal{H}$, which the answer only
reweights (prior $a_j w_{jk}$, with $a_j$ the firing's share of the error budget $A$), prices each
hypothesis by its own \protect\eprocess{}, estimates a one-period shift from the data (from
period~18, by replay of the frozen arm) and orders the $p/(p+h)$ quantile of the
posterior-predictive mixture. Its bound needs a known null and arrival law; as run, a null
simulation puts $\Expect_0[\Pi_\sigma]$ at up to $2.7A$ (Section~\ref{sec:limitations}). Each lane
ends in the episode from period~10 (every series is zero before it): commitment and cumulative
net reward against arm~1. Over all 32 lead-time episodes the hedge gains $+1{,}286$ per episode
on arm~1 (registered readout) and the gate, post hoc, $+414$.}
\label{fig:overview}
\end{figure*}

Language models are being proposed for replenishment decisions. On InventoryBench, 1{,}320 multi-period
lost-sales replenishment instances, the strongest published language-model strategy lets the
model name an order quantity every period and applies it untested~\cite{inventorybench}.
Whether that works cannot be read off the reported gain, because three decisions are welded
together: \emph{when} the model is consulted, \emph{what} it may say, and \emph{whether} its
answer is trusted. Reproducing the benchmark's operations research baseline also showed that
its published score belongs to a legacy simulator (Section~\ref{sec:background}), so the
evidence is partly artifactual as well as bundled.

We took the three decisions apart. A shared trigger decides when the model is asked; a closed
schema, the \shockspec{}, fixes what it may say (a typed claim about what changed in the
exogenous process, never an order); and an anytime-valid \eprocess{}, started strictly after
the proposal, decides whether the claim may touch an order. We preregistered the natural
version of the last step, a \emph{gate} that lets a hypothesis act only once its \eprocess{}
crosses the Ville threshold $1/\alpha_j$, and ran a pilot against a frozen decision rule.

The gate failed its own decision rule: the registered pilot read \emph{kill-or-reframe}, and
an exploratory rerun with live Gemini and Grok answers showed why. The gate rarely acted on a
false alarm (1 of 12 false-alert episodes with each model), but it acted on fewer than half of
the right-family proposals (43\% and 39\%), a median of 10 and 11 periods after they were made,
and then compiled whatever magnitude the model had stated, though no alert states a size
(Section~\ref{sec:results}).

This paper replaces the gate with \emph{certify-then-hedge}. The \eprocess{} no longer switches
the controller on; it prices a fixed set of action-aligned hypotheses. The language model's
answer only reweights that set (and fixes the onset window of the hypothesis it names), mixed
with a uniform prior so that a wrong answer has bounded cost. The size of a shock is estimated
from the data rather than taken from the model, and the order is the critical fractile of the
posterior-predictive mixture, so the controller moves toward a hypothesis as fast as the
evidence for it grows (Fig.~\ref{fig:overview}). Under a known null and arrival law, the construction keeps the
gate's error budget as a bound on expected exposure to a false hypothesis at every stopping
time, whatever the model says, and in the single-hypothesis case the gate's threshold is
analogous to a posterior-odds rule.


\subsection*{Contributions}

\begin{itemize}\itemsep0pt
\item \textbf{Certify-then-hedge}, a graded commitment rule for typed language-model
      hypotheses in a closed control loop, with an exposure bound that holds for any model
      output under a known null and arrival law, an odds-regret bound in cost units, and a
      simulation of how far the bound slips under the null we ran (Section~\ref{sec:hedge}).
\item \textbf{Two registered tests on fresh seeds}, exploratory with respect to the project's
      preregistration. At the middle cost ratio, against the gate and against the operations
      research baseline, with two proposer models, all four primary tests survive Holm
      correction and losses on no-shock episodes stay within the registered margin. At the
      benchmark's other two ratios the gains over the gate hold; at the higher one so does the
      gain over the baseline, and the model's answer adds value with one of the two models, but
      non-inferiority on no-shock episodes is not established there; on noisy lead times the
      gain over the baseline is not confirmed (Section~\ref{sec:results}).
\item \textbf{An account of where the gain comes from.} Against the baseline it comes almost
      entirely from lead-time shifts, and an alert-timed variant that ignores the model's
      answer is not distinguishable from it with the two primary models;
      the model's label lowers commitment on no-shock episodes and, post hoc, shows no
      detectable change in their reward. A model ladder (20 of 21 models run) shows, post hoc, that acting on proposals at once
      loses more the more readily a model commits, and that being right does not protect it.
\item \textbf{The negative pilot and the audit trail}: the gate's registered failure, every
      deviation and disclosure, and the audited benchmark divergence that motivated the
      controlled design.
\end{itemize}
